\chapter{User workflow}
\label{chap:ch5}

\par This chapter describes how the system is used by end users and explains typical workflows for common scenarios in daily use.

\section{Registration and login}

\par The first step for using the system is creating an account. The user must register by providing basic information such as name, email, and password. After registration, the account is created and can be used to access the system.

\par After registration, the user can log in using their email and password. Once logged in, the system creates a session for the user and allows access to the available features based on their role.

\section{Submitting a request}

\par After logging in, the user needs to complete their profile by adding personal information such as address and phone number.

\par After the profile is completed, the user can submit a request. The user can choose a request template from the available list and fill in the required information. Depending on the template, the user may also need to upload supporting documents.

\par Once the request is submitted, it is sent to the relevant department in the public institution. Employees review the submitted information and check whether all required documents are correct and complete.

\section{Review and feedback process}

\par During the review process, each submitted document is checked individually by the department employee. Each document can be either approved or rejected separately.

\par If a document is correct and meets the requirements, it is marked as approved. If a document is incorrect, missing information, or does not meet the requirements, it is rejected.

\par When a document is rejected, the employee can add a reason for the rejection. This helps the user understand what needs to be corrected or resubmitted. The user receives this feedback directly in the system.

\par After receiving feedback, the user can update the request by uploading corrected documents. The request is then reviewed again until all documents are approved.

\par When all documents in a request are approved, the request is marked as completed and the user is notified. The user can then collect the final document from the institution if needed.

\section{Administrator workflow}

\par Administrators are responsible for managing request templates in the system.

\par An administrator can create a new template for a specific type of request. Each template defines what documents are required and what information the user must provide. The administrator can also add example files for each required document to help users understand what is expected.

\par In addition, templates can be organized using tags. Tags help group similar templates together and make it easier to search and manage them inside the system.

\section{Department workflow}

\par Each request belongs to a specific department inside the institution. Employees from that department are responsible for handling incoming requests.

\par When a new request is submitted, it becomes available to department members. A department member can then claim a request, which means they become responsible for reviewing and processing it.

\par After claiming a request, other employees can see that it is already assigned and avoid duplicate work. The assigned employee continues the review process until the request is completed.

\par This system ensures that each request has a clear owner and improves organization inside the department.

\section{AI-assisted document processing}

\par The system includes AI-based tools to help with document processing. These tools are used to assist employees during the review of uploaded documents, especially when the document is unclear or hard to read.

\par One of the main use cases is when a document is blurry or of low quality. In this case, an AI service can analyse the image and estimate whether the document matches the expected type defined in the request template. This helps the employee decide faster if the document is valid or needs to be rejected.

\par Another feature is text extraction from images. If a document is scanned or uploaded as a photo, an OCR (Optical Character Recognition) service can extract the text content from it. This makes it easier to read and search through documents without manually inspecting every file.

\par In some cases, a text-to-speech (TTS) service can also be used to convert document content into audio. This can help users or employees who prefer listening instead of reading, or when reviewing long documents.

\par These AI tools do not make final decisions. They only provide suggestions or extracted information, while the final approval or rejection is always done by the assigned department employee.